%% ============================================================
%%  implementacion.tex  —  Sección 1.5
%% ============================================================

\subsection{Implementación del sistema}

\subsubsection{Arquitectura}

El sistema está organizado en módulos de cabecera (\texttt{.h}) independientes,
compilados junto con un programa principal (\texttt{.cpp}) que orquesta la
medición y el almacenamiento de resultados.
La Figura~\ref{fig:arquitectura} muestra la jerarquía de dependencias.

\begin{figure*}[t]
  \centering
  \begin{minipage}{0.85\linewidth}
  \footnotesize
  \begin{verbatim}
  graph.h          <- estructura Graph, Edge, generateRandomGraph(), INF
      |
      +-- dijkstra.h      <- DijkstraResult dijkstra(Graph, src)
      +-- bellman_ford.h  <- BellmanFordResult bellmanFord(Graph, src)
      +-- thorup.h        <- ThorupResult thorup(Graph, src, maxWeight)
      +-- det_mlogn.h     <- DetResult detMLogN(Graph, src, maxWeight)
                    |
                    +-- main_benchmark.cpp  (exps. 1-3, semilla 42)
                    +-- density_experiment.cpp (exp. 4, semillas 42/123/999)
  \end{verbatim}
  \end{minipage}
  \caption{Jerarquía de módulos del sistema de benchmark SSSP.}
  \label{fig:arquitectura}
\end{figure*}

\paragraph{Estructura de datos principal: \texttt{Graph}}
El grafo se representa mediante \textbf{listas de adyacencia}
(\texttt{std::vector<std::vector<Edge>{}>}), donde cada arista almacena el
vértice destino (\texttt{to}) y su peso (\texttt{weight}).
Este formato es eficiente para grafos dispersos y es la opción estándar para
los algoritmos de exploración iterativa evaluados.

\paragraph{Tipos de retorno}
Cada algoritmo devuelve una estructura específica que encapsula:
\begin{itemize}
  \item \texttt{dist}: vector de distancias mínimas desde el origen.
  \item \texttt{opsCount}: contador de relajaciones realizadas.
  \item Campos adicionales propios de cada algoritmo (e.g., \texttt{negativeCycle}
    en Bellman-Ford, \texttt{numComponents} en Thorup, \texttt{levels} en DMMSY).
\end{itemize}

\subsubsection{Tecnologías utilizadas}

\paragraph{C++17 — núcleo del benchmark}
Los cuatro algoritmos y los programas de benchmark están implementados en C++17,
compilados con \texttt{g++ -O2 -std=c++17 -Wall}.
Se utilizan exclusivamente componentes de la STL:
\texttt{priority\_queue}, \texttt{vector}, \texttt{deque}, \texttt{map},
\texttt{chrono} y \texttt{random}.

\paragraph{Python 3 — análisis y visualización}
El script \texttt{plot\_density.py} lee los archivos CSV generados en C++ y
produce cuatro gráficas PNG usando \texttt{matplotlib} (3.10.6),
\texttt{seaborn} (0.13.2) y \texttt{pandas} (2.3.2).

\paragraph{Simulador web — herramienta didáctica}
La carpeta \texttt{web/} contiene un simulador interactivo HTML/CSS/JS
(\texttt{index.html}, \texttt{app.js}, \texttt{style.css}) que:
\begin{itemize}
  \item Permite ajustar $n$ y $m$ mediante sliders en tiempo real.
  \item Calcula operaciones teóricas usando las fórmulas:
    $\text{Dijkstra} = (m+n)\log_2 n$,\;
    $\text{BF} = m \cdot n$,\;
    $\text{Det} = m\cdot(\log_2 n)^{2/3}$,\;
    $\text{Thorup} = m+n$.
  \item Anima la ejecución paso a paso sobre un grafo de demostración.
  \item Genera gráficas de complejidad con \texttt{Chart.js}.
\end{itemize}
Este simulador \textbf{no} realiza mediciones de tiempo real: sus operaciones
son conteos teóricos calculados en JavaScript monohilo.

\paragraph{Simplificaciones documentadas}

\begin{description}
  \item[\textbf{Thorup}] La implementación usa dos niveles de cubetas (fino y
    grueso) con ancho $\lfloor\sqrt{w_{\max}\cdot n}\rfloor$,
    en lugar de la jerarquía completa de componentes del paper.
    No implementa el modelo Word RAM ni opera sobre grafos no dirigidos
    (el paper requiere grafos no dirigidos).
    Complejidad resultante: $O(m + \sqrt{w_{\max}\cdot n})$.

  \item[\textbf{DMMSY}] La implementación calcula $K = \lceil\log^{2/3}n\rceil$
    niveles teóricos pero solo materializa el nivel~0 con ancho de cubeta
    $w_0 = \max\!\left(1,\lfloor w_{\max}\rfloor\right) = 1$.
    No existe promoción jerárquica entre niveles.
    El resultado es algorítmicamente equivalente a \textbf{Dial's Algorithm
    de un nivel}.
    El campo \texttt{levels} del resultado devuelve $K$ como valor informativo,
    no como niveles efectivamente utilizados.
\end{description}
